% ============================================================
\section{Introduction}
\label{sec:intro}
% ============================================================

The proliferation of Internet-of-Things (IoT) devices---projected to exceed
29 billion by 2030---has dramatically expanded the attack surface available
to adversaries.  Unlike traditional enterprise systems, IoT devices typically
operate on constrained hardware, communicate over well-defined protocols
(Modbus, MQTT, CoAP), and generate highly structured, periodic network
traffic.  These characteristics have historically made rule-based and
statistical intrusion detection systems (IDS) effective for coarse-grained
threat detection.

\textbf{The AI-powered threat gap.}
Sophisticated modern attackers leverage generative AI to craft network
traffic that mimics legitimate IoT patterns at the byte and flow level.
Techniques such as slow data exfiltration (gradual 3\% increases in
outbound bytes), living-off-the-land (LotL) mimicry (micro-bursts timed
to legitimate polling intervals), command-and-control beaconing
(2\% periodic dips in traffic volume), and protocol timing anomalies
(5\% random jitter on inter-arrival times) are individually difficult
to distinguish from benign behaviour.  We formalise these as
\emph{hard-negative} attacks: sequences where the statistical similarity
to benign traffic exceeds a stealth threshold $s \in \{0.85, 0.90, 0.95\}$
(see Section~\ref{sec:problem}).
Empirically, all five traditional IDS methods we evaluate---threshold,
signature, statistical, Isolation Forest, and weighted ensemble---achieve
less than 20\% detection rate on hard negatives at stealth level 95\%.

\textbf{Fighting AI with AI.}
We propose a fundamentally different defence strategy: \emph{use generative
AI to anticipate adversarial traffic}, and \emph{use foundation models to
detect it}.  Our system, \ours{} (\textbf{Diff}usion-based \textbf{IDS}),
proceeds in two stages.  First, we generate synthetic hard-negative attacks
using Diffusion-TS~\cite{zhang2024diffusionts}, a decomposition-aware
diffusion model, while ensuring that generated samples satisfy IoT protocol
constraints via a multi-tier validation framework.  Second, we fine-tune
Moirai~\cite{woo2024moirai}, a universal time-series forecasting foundation
model, using a combined negative log-likelihood (NLL) and supervised
contrastive loss~\cite{khosla2020supcon}.  The synthetic hard negatives
serve as difficult training examples that force the model to learn subtle
temporal signatures beyond what standard attacks provide.

\textbf{Contributions.}  We make four novel contributions:
\begin{enumerate}
  \item \textbf{Protocol-Constrained Hard-Negative Generation.}
    We extend Diffusion-TS with an IoT protocol validation layer
    (Modbus, MQTT, CoAP) that enforces packet-size, timing, and
    function-code constraints during generation, producing
    protocol-valid adversarial traffic at configurable stealth levels.

  \item \textbf{Foundation Model Adaptation for IDS.}
    We present the first systematic study of Moirai applied to network
    intrusion detection, including an engineering contribution: a fix for
    an in-place tensor operation in the \texttt{uni2ts} \texttt{PackedStdScaler}
    that silently broke gradient computation during fine-tuning.

  \item \textbf{Supervised Contrastive Fine-Tuning.}
    We introduce a combined loss $\mathcal{L} = \mathcal{L}_{\text{NLL}} +
    \lambda\,\mathcal{L}_{\text{SupCon}}$ ($\lambda=0.5$, $\tau=0.07$)
    with a projection head (384$\to$256$\to$128 dimensions) that
    bridges probabilistic forecasting and discriminative anomaly detection.

  \item \textbf{Stealth Attack Benchmark.}
    We contribute a reproducible benchmark: four IoT attack types with
    precise mathematical definitions at three stealth levels, evaluated
    against nine IDS methods on \dataset{}.
    Code, synthetic datasets, and a Streamlit demonstration are released
    at \url{https://github.com/anupamme/iotsf_demo}.
\end{enumerate}

\textbf{Paper organisation.}
Section~\ref{sec:related} reviews related work.
Section~\ref{sec:problem} formalises the problem.
Section~\ref{sec:method} describes \ours{}.
Section~\ref{sec:experiments} presents experimental results.
Section~\ref{sec:analysis} discusses findings and limitations.
Section~\ref{sec:conclusion} concludes.
